\originalchapter{图的语言与局部计数}\label{ch:basic}
\sourceof{6.042J 教材第11章; 18.315 图论部分的先修补全}
\originalsection{顶点、边与子图}
一张交通示意图可以变形而不改变道路连接关系. 图论首先保留这种关系, 再研究从关系本身能推出什么. 画在纸上的线段相交, 并不自动产生一个新顶点; 顶点必须在数学对象中指定.
\begin{definition}
有限简单无向图是 $G=(V,E)$, 其中 $V$ 为有限顶点集, $E$ 是 $V$ 的二元子集组成的集合. 边 $\{u,v\}$ 简记 $uv$, 其两个端点相邻. 简单图没有自环或平行边. 记 $n=|V|$, $m=|E|$, $N(v)=\{u:uv\in E\}$, $d(v)=|N(v)|$, 非空图的最小度与最大度分别为 $\delta(G)$、$\Delta(G)$. 空图约定最大度、顶点染色数及边染色数均为0, 最小度不定义.
\end{definition}
多重图允许两顶点间有若干条不同的边; 若还允许自环, 度数中每条自环计两次. 有向图的边为有序对, 称为弧; 顶点的入度和出度分别记 $d^-(v)$、$d^+(v)$. 除另行声明, 结构与染色章节均讨论有限简单无向图. 缩边、平面图对偶和流网络会显式允许平行边.

子图 $H=(W,F)$ 满足 $W\subseteq V$, $F\subseteq E$ 且各边端点属于 $W$. 诱导子图 $G[W]$ 保留 $W$ 内全部原边. 删除顶点 $v$ 时也删除与它关联的所有边; 删除边则不删端点. 补图 $\overline G$ 与 $G$ 顶点集相同, 两不同顶点在补图相邻当且仅当原图不相邻. 图同构是保持邻接关系的顶点双射, 因而编号与图形形状不是同构不变量.
\begin{example}
六个顶点组成一个六边形, 与两个互不相交的三角形, 是否同构?
\end{example}
\begin{solution}
两图都满足 $n=m=6$, 且每个顶点的度为2, 所以这些计数无法区分. 但前图只有一个连通分支, 后图有两个; 同构必须保持任意两顶点能否由路连接, 因而两图不同构. 度数列是有用的必要条件, 不是充分条件.
\end{solution}
完全图 $K_n$ 的任意两顶点相邻. 路图 $P_n$ 有 $n$ 个依次相连的顶点; 圈图 $C_n$ 在 $P_n$ 两端再加一条边, 其中 $n\ge3$. 完全二部图 $K_{a,b}$ 把顶点分为大小 $a,b$ 的两部分, 只在不同部分之间连边, 且所有这种边都存在.
\originalsection{握手引理与度数列}
\begin{theorem}[握手引理]\label{thm:handshake}
有限无向多重图满足 $\sum_{v\in V}d(v)=2m$. 有向图满足 $\sum_v d^+(v)=\sum_v d^-(v)=m$.
\end{theorem}
\begin{proof}
按关联的端点计数每条无向边, 非自环贡献两个不同端点, 自环在同一端点贡献两次. 有向弧恰有一个起点和一个终点, 对两种度分别求和均计每条弧一次.
\end{proof}
\begin{corollary}
任意有限无向图中, 奇度顶点的个数为偶数.
\end{corollary}
\begin{proof}
在度数总和中删去所有偶数项, 剩下若干奇数之和仍为偶数; 奇数项数只能为偶数.
\end{proof}
平均度是 $\overline d=2m/n$ ($n>0$). 它不意味着每个顶点都有相近的度, 但能保证某个顶点度不超过它、某个顶点度不小于它. 平面图的染色和极值图论都反复使用这个转换.
\begin{theorem}[Havel--Hakimi 归约]\label{thm:hh}
设 $d_1\ge d_2\ge\cdots\ge d_n\ge0$ 为整数, $d_1=k\le n-1$. 它是简单图的度数列, 当且仅当从 $d_2,\ldots,d_{k+1}$ 各减1, 保留其余项并重新排序所得序列是简单图的度数列. 归约出现负数时不可实现.
\end{theorem}
\begin{proof}
归约序列若可实现, 加一个新顶点并连接到被减1的 $k$ 个顶点即可. 反向需要说明为什么能指定最高度顶点的邻居. 在一张实现原序列的图中, 设最高度顶点为 $v$. 令 $T$ 为预定排序中的前 $k$ 个其他顶点. 若 $v$ 尚未邻接全部 $T$, 取 $x\in T\setminus N(v)$、$y\in N(v)\setminus T$, 于是 $d(x)\ge d(y)$. 我们把邻点 $y$ 换成 $x$.

具体说, 若找不到 $z\notin\{v,x,y\}$ 使 $xz\in E$ 而 $yz\notin E$, 则 $x$ 除可能的 $y$ 外的邻居都是 $y$ 的邻居; $y$ 还多出邻居 $v$, 导致 $d(y)>d(x)$, 矛盾. 取这样的 $z$, 删除 $vy,xz$, 加入 $vx,yz$. 新边原先不存在, 四个端点互异, 度数全不变. 每次交换使 $|N(v)\cap T|$ 增加1, 有限次后 $v$ 邻接指定的前 $k$ 个顶点. 删除 $v$ 得到归约序列.
\end{proof}
\begin{example}
判断 $(3,3,2,2,2)$ 与 $(3,3,3,1)$ 是否可实现.
\end{example}
\begin{solution}
第一列归约为 $(2,2,1,1)$, 再归约为 $(1,1,0)$, 最后为 $(0,0)$, 所以可实现. 第二列归约为 $(2,2,0)$, 再归约出现负数, 所以不可实现. 第二列的度数和虽为偶数, 但三个度3顶点都必须连接度1顶点, 已经矛盾. 握手条件本身不够.
\end{solution}
\originalsection{行走、路、圈与连通}
\begin{definition}
行走是相邻顶点组成的序列 $v_0,v_1,\ldots,v_\ell$, 长度为 $\ell$. 迹不重复边, 路不重复顶点. 闭行走满足 $v_0=v_\ell$; 圈除首尾重合外不重复顶点. 无向简单图的圈长度至少3. 两顶点之间有路时称它们连通; 极大的相互连通顶点集称为连通分支.
\end{definition}
从行走中删去两次出现同一顶点之间的闭段, 可得到同端点的路. 因而定义连通性时用行走或路等价. “有闭行走”与“有圈”则不同: 一条边来回走形成闭行走, 却没有圈. 若闭行走长度为奇数, 它一定含奇圈: 在最短奇闭行走中, 若有内部重复顶点, 分成的两个闭段之一仍为奇长, 与最短性矛盾.
\begin{theorem}\label{thm:bipartite}
无向图为二部图, 当且仅当它不含奇圈.
\end{theorem}
\begin{proof}
二部图的行走每一步都换到另一部分, 所以圈长度为偶数. 反向对每个分支选根 $r$, 按到根的最短路长度的奇偶分成两部分. 若某条边 $uv$ 的端点同属一部分, 从根到 $u$ 的最短路、边 $uv$ 和从 $v$ 回根的最短路组成奇闭行走, 因而含奇圈, 矛盾. 所有边都跨两部分.
\end{proof}
对连通图, 距离 $\dist(u,v)$ 是两顶点间最短路长度, 直径是所有距离的最大值. 圈图 $C_n$ 的直径为 $\lfloor n/2\rfloor$, 完全图 $K_n$ ($n\ge2$) 的直径为1. 非连通图可将不同分支之间的距离记为 $+\infty$; 讨论有限直径时必须先要求连通.
\originalsection{习题与解答}
\begin{exercise}\label{ex:basic1}
证明非空简单图中至少有两个顶点度数相同, 其中顶点数 $n\ge2$.
\end{exercise}
\begin{exercise}\label{ex:basic2}
证明 $n$ 个顶点、$m$ 条边的图有独立集大小至少 $n/(\Delta+1)$, 其中独立集指内部没有边的顶点集.
\end{exercise}
习题\ref{ex:basic1}的解答.
\begin{solution}
度数在 $0,1,\ldots,n-1$ 中. 度0和度 $n-1$ 不可能同时出现, 所以至多有 $n-1$ 个可取值, 对 $n$ 个顶点应用抽屉原理.
\end{solution}
习题\ref{ex:basic2}的解答.
\begin{solution}
不断取一个尚未删除的顶点放入独立集, 并删除它和所有邻居. 每次至多删除 $\Delta+1$ 个顶点; 选取的顶点彼此不邻接. 删除完全部顶点至少需 $n/(\Delta+1)$ 次. 整数意义下可向上取整.
\end{solution}
